% Source PDF pages 125--126; manual pages 2-96--2-97.
\subsection{Derivative Calculations}\label{sec:derivative-calculations}

The derivative calculation occurs inside the ``integrate all trajectories one time
step'' operation of \cref{fig:trajectory-calculation-flow}. A fourth-order
Runge--Kutta step calls the derivative function four times. The function
\texttt{compute\_derivs} in the \texttt{derivs} module performs each derivative
evaluation according to the sequence in
\cref{fig:derivative-calculation-flow}.

\begin{figure}[htbp]
  \centering
  \begin{tikzpicture}[
  >=Latex,
  node distance=0.28cm,
  flowbox/.style={draw,align=center,text width=5.3cm,minimum height=0.60cm,font=\scriptsize},
  guide/.style={draw,align=center,text width=5.05cm,minimum height=0.60cm,font=\scriptsize,fill=black!3},
  every path/.style={line width=0.6pt}
]
  \node[flowbox] (coords) {Compute coordinate-system unit vectors and the vehicle state in different coordinate systems};
  \node[flowbox,below=of coords] (rates) {Compute state rates};
  \node[flowbox,below=of rates] (gravity) {Compute acceleration of gravity};
  \node[flowbox,below=of gravity] (atmos) {Compute atmosphere properties};
  \node[flowbox,below=of atmos] (air) {Compute wind, airspeed, Mach number, dynamic pressure, and Reynolds number};
  \node[guide,below=of air] (att) {Compute body attitude, aerodynamic angles, and body unit vectors};
  \node[guide,below=of att] (forces) {Compute aerodynamic and propulsive forces};
  \node[guide,below=of forces] (derivs) {Compute accelerations and derivatives};
  \node[guide,below=of derivs] (loads) {Compute angular-acceleration rates and specific load factors};
  \node[flowbox,below=of loads] (rail) {Revise accelerations and derivatives for a rail launch or sled test};
  \node[flowbox,below=of rail] (other) {Compute other output variables};

  \foreach \a/\b in {coords/rates,rates/gravity,gravity/atmos,atmos/air,air/att,att/forces,forces/derivs,derivs/loads,loads/rail,rail/other}
    \draw[->] (\a)--(\b);

  \draw[->] (loads.east) -- ++(1.0,0) |- (att.east)
    node[pos=0.28,right,align=left,font=\scriptsize]
    {Loop to vary body attitude\\to satisfy guidance rules};
\end{tikzpicture}

  \caption{Derivative calculation flowchart.}
  \label{fig:derivative-calculation-flow}
\end{figure}

The derivative calculation is a function of time and the current vehicle state. TAOS
stores the propagated state in ECFC coordinates, so it first computes the vehicle
state in geocentric and geodetic coordinates
(\cref{sec:geocentric,sec:geodetic}). Many later calculations depend on those
coordinate representations.

TAOS next evaluates longitude, latitude, altitude, and related state rates from
\cref{sec:position-rates}. It then computes gravity
(\cref{sec:gravity-model}), atmospheric properties
(\cref{sec:atmosphere-model}), and the air-relative quantities such as airspeed,
Mach number, dynamic pressure, and Reynolds number
(\cref{sec:velocity-coordinate,sec:aerodynamic-forces}).

An initial body attitude starts the guidance loop. From that attitude, TAOS calculates
the body unit vectors and aerodynamic angles
(\cref{sec:body-fixed,sec:wind-coordinate}). The aerodynamic and propulsive
force tables are then evaluated, and the acceleration and state derivatives are formed
from the equations in \cref{sec:numerical-integration}.

Some guidance rules require geodetic acceleration components, flight-path-angle
rates, or specific load factors. TAOS evaluates those quantities near the end of each
guidance iteration and varies the body attitude until the commanded guidance
conditions are satisfied, as described in \cref{sec:guidance-rules}.

After the guidance loop converges, rail-launch or sled constraints may replace the
unconstrained acceleration according to \cref{sec:rail-launches}. The remaining
requested output variables are then calculated using the methods of
\cref{sec:output-variables}.
